% !TeX program = xelatex
\documentclass[margin = 5mm, tikz]{standalone}
\usepackage{xfrac}
%\usepackage{amsmath}
%\usepackage{unicode-math}
\usepackage[MnSymbol]{mathspec}
%\usepackage{mathspec}
%\setmainfont{Comic Sans MS}
%\setmathfont(Digits,Latin,Greek,Symbols)[Scale=MatchLowercase]{Comic Sans MS}
\setmainfont[Mapping=tex-text,Ligatures={NoRequired,NoCommon,NoContextual}]{Comic Sans MS}
\setmathfont(Digits,Latin,Greek)[Numbers={Lining,Proportional}]{Comic Sans MS}
%\setmathfont{Georgia}
%\setmathfont[range={"0000-"FFFF}]{Comic Sans MS}
%\setmathfont[range=up]{Comic Sans MS}
%\setmathfont[range=sfup]{Comic Sans MS}
%\setmathfont[range=it]{Comic Sans MS}
%\setmathfont[range=bfup]{Comic Sans MS}
%\setmathfont[range=bfit]{Comic Sans MS}
%\setsansfont{Comic Sans MS}  % Make \mathsf and \textsf also Calibri

%\usetikzlibrary{decorations.pathmorphing}
%\usetikzlibrary{decorations.pathreplacing}
\usetikzlibrary{shapes.geometric,decorations,decorations.pathmorphing}
\usetikzlibrary{intersections}
\usetikzlibrary{patterns}
\usetikzlibrary{calc}
\usetikzlibrary{shapes.misc}
\usetikzlibrary{spy}

\pgfdeclarelayer{bg}
\pgfdeclarelayer{bbg}
\pgfsetlayers{bbg, bg, main}

\definecolor{c1}{HTML}{ac0000}
\definecolor{c2}{HTML}{00a0d5}

\newcommand{\abs}[1]{{\left|{#1}\right|}}

\tikzset{point/.style={circle, draw=black, inner sep=2pt, fill=black}, empty point/.style={circle, draw=black, inner sep=2pt, fill=white}, cross/.style={cross out, draw, minimum size=2*(#1-\pgflinewidth), inner sep=0pt, outer sep=0pt},
	xkcd/.style = {decorate, decoration={random steps,segment length=0.8mm,amplitude=0.2pt}}}

\begin{document}
	\begin{tikzpicture}[spy using overlays={circle, magnification=4, connect spies, xkcd}]
		\pgfmathsetmacro{\xo}{1.7}
		\draw[thick, ->, xkcd] (-.5,0) -- (7,0);
		\draw[thick, ->, xkcd] (0,-1) -- (0,6);
		\draw[densely dashed] (\xo,-1) -- (\xo,6) node[pos=0, below]{$x_0$};
		\draw[thick, xkcd, opacity=1.0, domain=-0.5:{\xo-0.168}, samples=100] plot (\x,{-1/(\x-\xo)});
		\draw[thick, xkcd, opacity=1.0, domain={\xo+0.0248}:6.1, samples=100] plot ({\x+0.1},{0.8/(\x-\xo)*sin(5.4*\x r)+2.4}) node[right]{$f$};
		\draw[thick, xkcd, opacity=1.0, domain={\xo+0.168}:6.1, c2, samples=100] plot (\x,{1/(\x-\xo)}) node[right]{$g$};
		\begin{scope}
			\clip (-.5,-1) rectangle (7,5.95);
			\draw[thick, xkcd, opacity=1.0, domain=-0.5:{\xo-0.009}, c2, samples=100] plot ({\x-0.06},{-((\x-1.5)^2*(\x-0.1)*(\x-0.7))/(\x-\xo)+0.3});
		\end{scope}
%		\draw[thick, c2, xkcd, domain=2:6.3] plot (\x,{0.88*(\x-4.4) + 0.36})  node[above]{$\epsilon_g$};
%		\draw[thick, xkcd, domain=2:6.3] plot (\x,{-0.30*(\x-4.4) - 0.14})  node[right]{$\epsilon_f$};
%		\node[circle, fill=c1!10, draw=c1, inner sep=1.5pt] (a) at (4,0){};
%		\node[label={right, yshift=0.2cm}:{$x_1$}] (b) at (4.4,4){};
%		\node[label={below}:{$x_0$}] (c) at (3.6,3.8){};
%		\draw[c1, xkcd, thick, densely dashed] ($(-1,4.8)!(b)!(7,4.8)$) -- ($(-1,3.6)!(b)!(7,3.6)$);
%		\spy[c1, size=6cm, xkcd] on (a) in node [c1, opacity=0.08, xkcd] at (3.5,4);
	\end{tikzpicture}
\end{document}